\SourceLeaf{401}%
or
\[
\alpha=1,\quad\beta=0
\]
\par
The denominator in the expressions for \(\xi,\eta\) is twice the area of the triangle.

\par\medskip\Needspace{5\baselineskip}
\pdfbookmark[2]{II.}{carnot-zusatz-ii}
\centerline{\textbf{II.}}\nobreak\smallskip
There is a very simple way to show that the perpendiculars from the vertices of a triangle to the opposite sides meet at one point.

\begingroup\rightskip=0pt plus 1em
Let the given triangle be \(BDF\) and the perpendiculars just mentioned be \(\overline{BI},\overline{DG},\overline{FH}\).\par
\endgroup

Through each vertex of the triangle draw a line parallel to the opposite side. These lines meet at the points \(A,C,E\). Thus \(\overline{FH}\) is also perpendicular to \(\overline{AE}\), \(\overline{GD}\) to \(\overline{CE}\), and \(\overline{BI}\) to \(\overline{AC}\). Moreover,
\[
\begin{aligned}\overline{AB}&=\overline{BC}\\\overline{ED}&=\overline{DC}\\\overline{AF}&=\overline{FE}\end{aligned}
\]
Draw a circle around the triangle \(ACE\). Its centre lies on each of \(\overline{BI}\), \(\overline{DG}\) and \(\overline{FH}\), so the three lines must meet at one point.

\textsc{Puissant}, in his Recueil des propositions de Géométrie, gives an elegant analytical proof and adds a geometrical one that does not have the same merit.

\par\medskip\Needspace{5\baselineskip}
\pdfbookmark[2]{First handwritten remark.}{carnot-handschriftlich-i}
\begin{center}[\textit{First handwritten remark.}]\end{center}
\hspace*{\parindent}Let \(\alpha,\beta,\gamma\) be the complex numbers corresponding in their natural order to the three vertices of a triangle, let \(A,B,C\) be its three angles, and let \(u\) be the complex number for the centre of the [circumscribed] circle. Then
\begingroup\footnotesize
\[
\begin{aligned}
2u&=\alpha+\beta+(\beta-\alpha)\operatorname{cotg}C.i=\alpha(1-i\operatorname{cotg}C)+\beta(1+i\operatorname{cotg}C)\\
&\hphantom{=\alpha+\beta+(\beta-\alpha)\operatorname{cotg}C.i}=\beta(1-i\operatorname{cotg}A)+\gamma(1+i\operatorname{cotg}A)\\
&\hphantom{=\alpha+\beta+(\beta-\alpha)\operatorname{cotg}C.i}=\gamma(1-i\operatorname{cotg}B)+\alpha(1+i\operatorname{cotg}B)
\end{aligned}
\]
\endgroup
